\section{总结与延伸}

本节把整期压缩成一个学习框架。第一，模型范式不是孤立论文，而是硬件、数据、架构和训练技巧对齐后的结果。第二，Infra 与数据是模型能力的放大器，没有显存分片、集群效率和高质量数据，Scaling Law 只是纸面规律。第三，语言模型的发展是从表示到预训练，再到后训练和对齐的过程。第四，多模态的发展是视觉理解、图像生成、图文对齐和 Transformer 化逐步合流的过程。第五，读论文的关键不是读完，而是建立边界意识。

\begin{importantbox}{把 EP117 放进张小珺 AI 队列}
EP119 讲 Attention 架构考古，EP118 讲 CEO 如何理解 VLA、Agent OS 和平台，EP117 则把这些问题往历史深处拉：为什么 Transformer 会赢，为什么 Infra 和数据成为主线，为什么语言模型和多模态最终会汇合。
\end{importantbox}

\subsection{36 篇论文的压缩表}

\begin{longtable}{p{0.22\textwidth}p{0.30\textwidth}p{0.38\textwidth}}
\toprule
主线 & 论文/节点 & 教学意义 \\
\midrule
模型范式 & Brook、AlexNet、ResNet & GPU、深度学习和残差网络共同改变可训练边界。 \\
模型范式 & seq2seq、Attention、Transformer & 序列建模从压缩状态走向直接建模 token 关系。 \\
模型范式 & AlphaGo Zero、MoE、CoT、LoRA、ReAct & 强化学习、专家路由、推理提示、低成本适配和 Agent 形成后续系统能力。 \\
模型范式 & The Bitter Lesson & 长期看，通用计算和规模化学习反复压过手工结构。 \\
Infra/数据 & ZeRO、MegaScale & 分片和集群工程决定大模型能否被训练。 \\
Infra/数据 & Scaling Law、Chinchilla & 参数、数据和计算预算需要配比。 \\
Infra/数据 & LAION-5B、RefinedWeb & 开源数据和网页数据工程决定可训练语料边界。 \\
语言模型 & Word2Vec、Google Translate & 词表示和线上神经模型部署是语言模型前史。 \\
语言模型 & GPT-1、BERT、GPT-2、GPT-3 & 预训练范式从理解任务走向生成和规模涌现。 \\
语言模型 & InstructGPT、Tulu 3 & 后训练、指令对齐和开源复现让模型可用。 \\
多模态 & DeepVideo、双流网络 & 视频理解把深度学习从静态图像推向时间维度。 \\
多模态 & GAN、Diffusion、DDPM & 生成模型从对抗训练走向扩散稳定生成。 \\
多模态 & ViT、CLIP、Stable Diffusion、DiT & 视觉 Transformer、图文对齐和扩散 Transformer 走向融合。 \\
\bottomrule
\end{longtable}

\subsection{拓展阅读}

\begin{itemize}
\item 对模型范式感兴趣，可以继续对照 EP119 的 Attention 架构综述，理解 Transformer 之后为什么仍然要继续“雕 Attention”。
\item 对物理世界 AI 感兴趣，可以继续对照 EP118、EP120、EP121，把 VLA、世界模型和机器人数据放进同一张图。
\item 对 Agent 感兴趣，可以继续对照 EP139 的 Agent 技术史和 EP138 的后训练/Agent 范式，理解 CoT/ReAct 如何变成真实系统。
\end{itemize}

\end{document}
